\chapter{Rules}

본 장에서는 2026 과학퀴즈 공식 게임 기획안과 경기 규정에 따라 RPSC의 규칙을 정리한다.

\section{Board and Pieces}

RPSC는 $8\times8$ 체스판에서 진행한다. 각 팀은 정육면체 주사위 말 네 개를 사용한다. 각 말의 여섯 면에는 가위, 바위, 보가 두 면씩 그려져 있으며, 같은 손 모양의 두 면은 서로 마주본다. 초기 배치와 말의 방향은 고정되어 있고, 시작할 때 각 면의 손목은 해당 팀의 몸 쪽을 향한다. 공식 경기에서 \POSTECH{}은 빨간색, \KAIST{}는 파란색 말을 사용한다.

\begin{figure}[htbp]
    \centering
    \begin{tikzpicture}[x=0.95cm,y=0.95cm]
        \foreach \x/\y in {0/0,1/0,2/0,3/0,1/1,1/-1} {
            \draw[rpsc board line,fill=white] (\x,\y) rectangle ++(1,1);
            \draw[rpsc board border] (\x,\y) rectangle ++(1,1);
        }
        \rpscglyph{0.50}{0.56}{P}{0}{black}
        \rpscglyph{1.50}{0.56}{S}{0}{black}
        \rpscglyph{2.50}{0.56}{P}{0}{black}
        \rpscglyph{3.50}{0.56}{S}{0}{black}
        \rpscglyph{1.50}{1.56}{R}{0}{black}
        \rpscglyph{1.50}{-0.44}{R}{0}{black}
    \end{tikzpicture}
    \caption{주사위 말의 전개도. 같은 손 모양의 두 면은 서로 마주본다.}
    \label{fig:dice-net}
\end{figure}


이 책에서는 편의상 가위, 바위, 보를 각각 S, R, P로 적는다. 말의 기본 이동거리는 이동을 시작할 때의 윗면에 따라 정한다.

\begin{center}
    \renewcommand{\arraystretch}{1.15}
    \begin{tabular}{@{}ccc@{}}
        \toprule
        윗면 & 표기 & 기본 이동거리 \\
        \midrule
        가위 & S & 3칸 \\
        바위 & R & 4칸 \\
        보 & P & 5칸 \\
        \bottomrule
    \end{tabular}
\end{center}

\section{Movement and Combat}

한 번의 이동에서는 자신의 말 하나를 선택하여 정해진 칸 수만큼 한 칸씩 \emph{굴린다}. 매 칸마다 상하좌우 가운데 한 방향을 선택할 수 있고, 한 번의 이동 안에서 여러 방향을 자유롭게 조합할 수 있다. 다만 바로 직전에 있던 칸으로 즉시 되돌아갈 수는 없다. 다른 말이 놓인 칸에는 들어갈 수 없으며, 그 칸을 통과할 수도 없다.

말을 한 칸 굴릴 때마다 정육면체의 위치와 방향이 함께 바뀌므로 윗면도 바뀔 수 있다. 한 번의 이동에서 굴려야 하는 전체 칸 수는 이동을 시작할 때 정하며, 이동 도중 윗면이 바뀌더라도 다시 계산하지 않는다.

이동을 모두 마친 뒤, 이동한 말이 상대 말과 상하좌우로 인접하면 가위바위보 대결을 한다. 일반적인 가위바위보 규칙에 따라 이긴 말은 남고 진 말은 제거된다. 두 말의 손 모양이 같으면 어느 말도 제거되지 않는다. 이동을 마쳤을 때 두 개 이상의 상대 말과 동시에 대결이 성립하는 칸에는 멈출 수 없지만, 그러한 칸을 이동 도중 지나가는 것은 가능하다.

\section{Quiz and Items}

본 경기에서는 사전에 검수된 과학퀴즈 본 문제 20문제와 RPSC를 번갈아 진행한다. 각 문제의 풀이 시간은 1분이며, 두 팀의 정오 여부에 따라 다음 행동이 결정된다.

\begin{center}
    \renewcommand{\arraystretch}{1.18}
    \begin{tabular}{@{}L{0.24\textwidth}L{0.28\textwidth}L{0.34\textwidth}@{}}
        \toprule
        퀴즈 결과 & 보드 진행 & 아이템 \\
        \midrule
        양 팀 모두 정답 & 선공, 후공 순서로 각각 이동 & 없음 \\
        양 팀 모두 오답 & 선공, 후공 순서로 각각 이동 & 없음 \\
        한 팀만 정답 & 보드 이동 없음 & 정답 팀이 원하는 아이템 1개 획득 \\
        \bottomrule
    \end{tabular}
\end{center}

두 팀의 정오 여부가 같아 보드 이동이 발생하면 선공 팀이 먼저, 후공 팀이 다음으로 움직인다. 각 팀에는 30초가 주어진다. 자신의 말 하나를 선택하여 기본 이동 규칙에 따라 이동하거나, 보유한 아이템 하나를 적용한 뒤 이동할 수 있다. 한 번의 이동에 사용할 수 있는 아이템은 최대 한 개이며, 사용한 아이템은 상대 팀에게 공개한다. 제한 시간 안에 행동하지 못하면 아이템을 사용하지 않고 무작위로 말을 이동한다.

한 팀만 정답을 맞히면 그 팀은 다음 세 종류 가운데 원하는 아이템 하나를 획득하며, 이때 보드 이동은 하지 않는다. 아이템 선택 시간은 30초이고, 선택한 종류는 상대 팀에게 공개한다. 제한 시간 안에 선택하지 못하면 무작위로 아이템 하나를 획득한다. 아이템은 개수 제한 없이 보유할 수 있으며, 이후 자신의 이동 차례에 사용할 수 있다.

\begin{center}
    \renewcommand{\arraystretch}{1.18}
    \begin{tabular}{@{}L{0.20\textwidth}L{0.68\textwidth}@{}}
        \toprule
        아이템 & 효과 \\
        \midrule
        Push & 말을 상하좌우 가운데 한 방향으로 한 칸 민다. 이때 말을 굴리지 않으므로 윗면과 방향은 변하지 않는다. 그 위치에서 원래의 기본 이동거리만큼 굴린다. \\
        Rotation & 말을 현재 칸에 세운 채 보드에 수직인 축을 중심으로 $90^\circ$ 회전시킨다. 위에서 내려다본 기준으로 시계방향 또는 반시계방향을 선택한다. 이 동작은 말을 굴리는 것이 아니므로 윗면은 바뀌지 않고, 따라서 기본 이동거리도 바뀌지 않는다. 회전으로 바뀐 방향을 유지한 채 원래의 기본 이동거리만큼 굴린다. \\
        Step & 기본 이동거리보다 한 칸 더 가거나 한 칸 덜 간다. 이 책의 기보에서는 이를 \texttt{StL}, \texttt{StS}로 적는다. \\
        \bottomrule
    \end{tabular}
\end{center}

이 책의 기보에서는 Rotation의 반시계방향 회전을 \texttt{RoL}, 시계방향 회전을 \texttt{RoR}로 적는다. Rotation은 \emph{제자리 회전}이고, 이후의 이동은 한 칸씩 주사위를 \emph{굴리는 동작}이라는 점을 구분한다.

공식 규칙에서는 직전에 마지막으로 말을 움직인 팀이 다음 보드 진행의 후공이 되고 상대 팀이 선공이 된다. 한 번의 보드 진행에서 두 팀이 선공, 후공 순서로 한 번씩 움직이므로 최초 선후공이 정해진 뒤에는 같은 순서가 유지된다. 과학퀴즈가 끝나기 전에 한 팀의 말 네 개가 모두 제거되면, 양 팀의 말 여덟 개를 모두 초기 배치와 방향으로 되돌린다. 이때 이미 얻은 점수와 보유 아이템, 선후공은 그대로 유지한다.

\section{Scoring and Result}

과학퀴즈 한 문제에 정답하면 1점을 얻고, 가위바위보 대결에서 상대 말을 제거하면 2점을 얻는다. 경기 점수는 두 점수의 합으로 계산한다.

본 문제 진행이 끝나면 총점이 높은 팀이 승리한다. 총점이 같으면 과학퀴즈를 더 많이 맞힌 팀이 승리한다. 과학퀴즈 정답 수도 같으면 RPSC는 더 진행하지 않고 연장전을 치른다. 연장전은 수학, 물리학, 화학, 생명과학, 컴퓨터공학에서 한 문제씩 출제된 최대 5문제를 무작위 순서로 풀며, 한 팀이 앞서는 순간 즉시 종료하여 그 팀을 승자로 한다.
